\makeproblem{DAOVANG}{ĐÀO VÀNG}{1.0s}{1024 MB}

Đào vàng là trò chơi khá phổ biến và có nhiều phiên bản. Hãy cùng xét một phiên bản của trò chơi này. Có $N$ thỏi vàng được cố định ở các vị trí $X_1, X_2, X_3, \dots, X_N$ trên một trục nằm ngang. Nếu người chơi đào ở vị trí $X$ với máy khoan có lực đập $R$ thì có thể lấy được các thỏi vàng cách vị trí $X$ tối đa $R$ đơn vị chiều dài hay các thỏi vàng có vị trí nằm trong khoảng $[X - R; X + R]$. Người chơi được đào tối đa $K$ lần và lực đập $R$ là giống nhau ở các lần đào. Nếu người chơi chọn lực đập $R$ càng nhỏ thì số điểm đạt được càng cao và ngược lại. Người chơi được thực hiện tối đa $K$ lần đào, hãy giúp người chơi chọn lực đập $R$ nhỏ nhất để có thể đào hết $N$ thỏi vàng.

\medskip
\textbf{Yêu cầu:} Cho trước vị trí của $N$ thỏi vàng, hãy viết chương trình tìm giá trị nguyên $R$ bé nhất sao cho người chơi có thể lấy được $N$ thỏi vàng sau tối đa $K$ lần đào.

\makesection{Input}
\begin{itemize}
    \item Dòng đầu chứa hai số nguyên $N$ và $K$ lần lượt cho biết số lượng thỏi vàng và số lần đào tối đa.
    \item Dòng thứ $i$ trong $N$ dòng tiếp theo cho biết vị trí $X_i$ ($0 \le X_i \le 10^9$) của thỏi vàng thứ $i$.
\end{itemize}

\makesection{Output}
\begin{itemize}
    \item Ghi ra một số nguyên là giá trị lực đập $R$ bé nhất để lấy được $N$ thỏi vàng sau tối đa $K$ lần đào.
\end{itemize}

\begin{makesubtasks}
    \subtask{20} $K = 1$ và $N \le 1000$.
    \subtask{20} $K = 2$ và $N \le 10\,000$.
    \subtask{60} $K \le 20$ và $N \le 50\,000$.
\end{makesubtasks}

\makesection{Ví dụ}

\subsubsection*{Ví dụ 1}
\begin{makesample}
\begin{lstlisting}
6 1
2
20
6
5
4
17
\end{lstlisting}
\nextsample
\begin{lstlisting}
9
\end{lstlisting}
\end{makesample}

\subsubsection*{Ví dụ 2}
\begin{makesample}
\begin{lstlisting}
6 2
2
20
6
5
4
17
\end{lstlisting}
\nextsample
\begin{lstlisting}
2
\end{lstlisting}
\end{makesample}